% Felipe Muñoz Casasola
% This is free and unencumbered software released into the public domain.

% Anyone is free to copy, modify, publish, use, compile, sell, or
% distribute this software, either in source code form or as a compiled
% binary, for any purpose, commercial or non-commercial, and by any
% means.

% In jurisdictions that recognize copyright laws, the author or authors
% of this software dedicate any and all copyright interest in the
% software to the public domain. We make this dedication for the benefit
% of the public at large and to the detriment of our heirs and
% successors. We intend this dedication to be an overt act of
% relinquishment in perpetuity of all present and future rights to this
% software under copyright law.

% THE SOFTWARE IS PROVIDED "AS IS", WITHOUT WARRANTY OF ANY KIND,
% EXPRESS OR IMPLIED, INCLUDING BUT NOT LIMITED TO THE WARRANTIES OF
% MERCHANTABILITY, FITNESS FOR A PARTICULAR PURPOSE AND NONINFRINGEMENT.
% IN NO EVENT SHALL THE AUTHORS BE LIABLE FOR ANY CLAIM, DAMAGES OR
% OTHER LIABILITY, WHETHER IN AN ACTION OF CONTRACT, TORT OR OTHERWISE,
% ARISING FROM, OUT OF OR IN CONNECTION WITH THE SOFTWARE OR THE USE OR
% OTHER DEALINGS IN THE SOFTWARE.

% For more information, please refer to <https://unlicense.org/>

\documentclass{article}

\usepackage[spanish]{babel}
\usepackage[margin=3cm]{geometry}
\usepackage{microtype}
\usepackage[pdfusetitle]{hyperref}
\usepackage{minted}

\title{Reinstalación completa del sistema paso por paso}
\author{Felipe Muñoz Casasola}
\date{24 de enero de 2026; revisado el 21 de febrero de 2026}

\begin{document}
    \maketitle
    \clearpage
    \tableofcontents
    \clearpage
    \section{Copia de seguridad previa}
    Seguir las indicaciones de los documentos de copias de seguridad para hacer
    una copia de seguridad del sistema antes de reinstalar. Recuerda los
    marcadores del navegador, los perfiles (cerrar las cuentas si no se pueden
    trasladar al nuevo sistema).

    \section{Documentación a tener preparada}
    Este documento PDF y todos los documentos PDF de gestión del sistema, el
    manual de instalación de Debian.

    \section{Instalación básica}
    Uso la instalación por red mínima, que tiene el tamaño de imagen más
    pequeño. El modo de instalación es el modo gráfico experto. Selecciono que
    se separen las particiones \texttt{home},
    \texttt{tmp} y \texttt{var}. El entorno de escritorio a instalar (por si
    fallara el dwm) es LXQt. GRUB está instalado con la opción para detectar
    otros sistemas operativos, solamente por si acaso hiciera falta. El
    superusuario (raíz) está activado, y tiene la misma contraseña que el único
    usuario que se crea, el usuario usuario.

    \subsection{Vim}
    Instalar vim provisionalmente (tendré que recompilarlo para tener el Vimtex
    posteriormente) y eliminar \texttt{nano} y \texttt{FeatherPad} (o el editor
    de texto que venga por defecto con el entorno de escritorio.

    \subsection{Añadir \texttt{sbin} a la ruta \textit{PATH}}
    Añadir en el \textit{PATH} que se encuentra en \texttt{/etc/profile} las
    rutas de \texttt{sbin} y \texttt{/usr/local/sbin} en los dos casos del
    \texttt{if} en el que se comprueba si el usuario es el superusuario o no.

    \subsection{Permitir al usuario usar \texttt{sudo}}
    Para ello tengo que agregar al usuario al grupo de usuarios
    \texttt{sudo}. Por ejemplo con el comando:

    \texttt{usermod -a -G sudo usuario}

    Y para que el cambio tenga efecto se tiene que reiniciar el sistema.

    \section{Arreglar \textit{firmware}}
    A pesar de haber seleccionado previamente durante la instalación la
    inclusión del \textit{firmware} no libre, es muy posible que tenga
    problemas. Principalmente suelo tener con el WiFi y con la GPU. Para
    arreglarlos está la herramienta \texttt{isenkram} y es de gran
    utilidad la wiki de Debian, en concreto \href{https://wiki.debian.org/Firmware}
    {https://wiki.debian.org/Firmware}.

    No obstante, es posible que algunos de los archivos de \textit{firmware}
    que faltan no se encuentren de la anterior forma, y por tanto sea
    necesario buscarlos en el kernel. Esto está sacado de la página
    anteriormente referenciada, que se encuentra archivada
    \href{https://web.archive.org/web/20260131222721/https://wiki.debian.org/Firmware}
    {aquí}.

    Por tanto, ya en los archivos de \textit{firmware} del kernel debería
    poder encontrarlos. Aunque no siempre es posible encontrarlos todos.

    En concreto se trata del repositorio de \textit{firmware} del kernel,
    no en el de el propio kernel. A día de hoy dicho repositorio se encuentra
    \href{https://git.kernel.org/pub/scm/linux/kernel/git/firmware/linux-firmware.git/tree}
    {aquí}.

    \section{Archivos de configuración a guardar}
    Los archivos de configuración se deben guardar en un dispositivo con un
    sistema de archivos que permita los enlaces simbólicos, pues el tema del
    ratón requiere guardar algunos enlaces simbólicos.

    \begin{enumerate}
    	   \item QEMU:
    	           \subitem La carpeta «Máquinas virtuales» dentro de la carpeta
		   raíz del usuario, que tiene los discos duros de las máquinas
		   virtuales que tengo.
    	   \item Navegador Mullvad
    	   	   \subitem Archivo de marcadores.
    	   	   \subitem Archivo de configuración de NoScript.
    	   	   \subitem Archivo de configuración de Ublock Origin.
    	   	   \subitem Archivo de configuración de Privacy Badger.
    	   	   \subitem Archivo de configuración de Dark Reader.
    	   	   \subitem Archivo de configuración de Canvas Blocker.
               \subitem Certificados de usuario.
    	   \item SearxNG
    	   	   \subitem Archivo de configuración en
		   \texttt{$\sim$/.config/searxng/settings.yml}.
		   \subitem \textit{Quadlet} en
		   \texttt{$\sim$/.config/containers/systemd/searxng.container}.
    	   \item Tema del ratón
    	           \subitem Carpeta
		   \texttt{/usr/share/icons/material\_light\_cursors}.
	   \item Synchting
		   \subitem Archivos (no directorios) en la carpeta
		   \texttt{$\sim$/.local/stat/syncthing/}.
		   \subitem Carpetas compartidas en Syncthing, pues si no
		   existen con la carpeta \texttt{.synchting} dentro, Syncthing
		   se bloquea.
	   \item Rclone
		   \subitem Archivo de configuración de Rclone, que actualmente
		   se encuentra en
		   \texttt{$\sim$/.config/r\-clone/\-rclone.conf}.
		   \subitem Servicio para que automáticamente (activarlo para
		   que se ejecute en cada arranque) se guarde una copia de
		   seguridad cifrada en MEGA, actualmente en
		   \texttt{$\sim$/.config/systemd/user.control/rclone.service}. Debo tener la
		   carpeta Documentos previamente restaurada para que al
		   reiniciar el sistema no borre la copia de seguridad en la
		   nube porque los documentos estén vacíos.
		   \subitem Otro servicio pero para hacer copia de seguridad del
		   escritorio, actualmente en
		   \texttt{$\sim$/.config/systemd/user.control/rclone\_e.service}. Tengo que
		   haber restaurado previamente la carpeta Escritorio.
	   \item \LaTeX{}
		   \subitem Carpeta \texttt{$\sim$/texmf}, donde se guardan los
		   paquetes de \LaTeX{} que haya creado.
	   \item Vim
		   \subitem Carpeta \texttt{.vim} en el directorio raíz del
		   usuario.
		   \subitem Archivo \texttt{.vimrc} en la carpeta raíz del
		   usuario.
	   \item KeePassXC
		   \subitem Tengo que asegurarme de que los archivos de las
		   contraseñas los haya guardado.
           \subitem También tengo que guardar la lista de frases contraseña
           que tengo en español. De la que guardo una copia en este mismo
           USB.
           \subitem Y la base de datos de los códigos TOTP que tengo en la carpeta
           Compartir con nombre de archivo «códigos.kdbx».
	   \item DWM
		   \subitem Carpeta \texttt{.barra} en el directorio raíz del
		   usuario.
		   \subitem Archivo de configuración de DWM
		   \texttt{config.h}.
		   \subitem Carpeta \texttt{.fonts} en la carpeta raíz del
		   usuario.
		   \subitem Carpeta \texttt{dwm} donde se encuentra el archivo
		   \texttt{xinitrc}, que se encuentra en la misma carpeta en la
		   que el gestor de escritorio LXQt esté instalado por defecto,
		   que actualmente es en \texttt{/etc/xdg}.
		   \subitem Archivo \texttt{dwm.desktop} para que aparezca como
		   un entorno de escritorio a seleccionar cuando se inicia
		   sesión, actualmente se encuentra en
		   \texttt{/usr/share/xsessions/dwm.desktop}.
	   \item Terminal
		   \subitem Archivo \texttt{.bashrc}.
        \item Lf
            \subitem Archivo \texttt{lfrc}.
            \subitem Archivo \texttt{icons}.
            \subitem Archivo \texttt{colors}.
        \item Forgejo
            \subitem Archivo \texttt{/etc/containers/systemd/forgejo.container}.
            \subitem Archivo \texttt{/etc/containers/systemd/forgejo.network}.
            \subitem Archivo \texttt{/etc/containers/systemd/forgejo-data.volume}.
            \subitem Directorio
	    \texttt{/var/lib/containers/storage/volumes/forgejo-data/\_data/gitea}.
            \subitem Directorio \texttt{$\sim$/.ssh}
    \end{enumerate}

    \section{Instalación del navegador Mullvad}
    Accedo al navegador Firefox que viene preinstalado, y me meto en la página
    oficial del navegador para las instrucciones de cómo añadir el repositorio
    del navegador Mullvad a los repositorios del sistema para que se pueda
    actualizar automáticamente con la orden de actualizar.

    Una vez instalado el navegador Mullvad, hay que eliminar el navegador
    Firefox.

    También tengo que instalar los siguientes complementos (y cargar la
    configuración que guardé previamente sobre ellos):
    \begin{itemize}
    	   \item uBlock Origin.
    	   \item NoScript.
    	   \item Privacy Badger.
    	   \item CanvasBlocker.
    	   \item Dark Reader.
    \end{itemize}
    
    Además, tengo que instalar la extensión Decentraleyes, de la que no guardo
    configuración. Y eliminar el complemento del navegador Mullvad que viene por
    defecto.

    \section{Instalación de SearxNG}
    Si todavía existe la opción, es mejor instalarlo con el contenedor,
    preferiblemente \texttt{podman}, pues me suele dar fallos al instalar con
    las otras alternativas que se ofrecen actualmente. Una vez hecho esto usar
    la guía de la página de SearxNG para redirigir la configuración que guardaré
    en \texttt{$\sim$/.config/searxng/} a la configuración del contenedor. Para
    generar la
    configuración simplemente descargaré el archivo de configuración del
    repositorio de GitHub y lo compararé con el que he guardado de
    configuración, para activar y desactivar las mismas opciones. Esto lo hago
    así porque es posible que algún motor de búsqueda haya quedado obsoleto y al
    aparecer en el antiguo archivo de configuración el servidor falle.

    Para más información consultar el manual de SearxNG.

    \section{Instalación de QEMU}
    Para volver a instalar QEMU en el nuevo sistema tengo que seguir los
    siguientes pasos:

    \begin{enumerate}
    	   \item \texttt{sudo apt install virt-manager}. Esto instala el gestor
		   de las máquinas virtuales junto con QEMU.
    \end{enumerate}
    
    Una vez hecho esto ya tengo instalado el sistema base. Solamente queda
    copiar los discos duros virtuales de la configuración que se guardó, y
    configurarlas. Para ello simplemente tengo que crear nuevas máquinas
    virtuales usando la opción para importar con una imagen de disco existente,
    que previamente guardé antes de reinstalar el sistema.

    Después de haber creado las máquinas virtuales tengo que ejecutar el comando
    \texttt{sudo virsh net-autostart default}, gracias al usuario Querzion en
    \href{https://forum.level1techs.com/t/solved-qemu-network-default-is-not-active-booting-from-hard-disk/194698}
    {level1techs}

    \section{Tema del ratón}
    Instalar el tema \textit{Material Cursors} que se encuentra en la página
    \href{https://www.gnome-look.org}{www.gnome-look.org}. Que se encuentra en
    GitHub en la página
    \href{https://github.com/varlesh/material-cursors}
    {https://github.com/varlesh/material-cursors}. Desinstalar las herramientas
    para instalarlo inmediatamente después de la instalación.

    Nótese que el guion \texttt{build.sh} no me funciona, tengo que usar el que
    he guardado corregido. Actualmente tarda mucho en terminar de «compilar» los
    cursores.

    Si no funciona el instalador del tema del ratón, se puede hacer manualmente,
    para ello se tiene que actualizar la alternativa actualmente llamada
    \texttt{x-cursor-theme}. Para ello, veo a qué archivo apunta por defecto,
    para quedarme con la ruta en la que se encuentran los temas del ratón. Y
    entonces pego en la carpeta padre (donde se estarán almacenando todos los
    temas del ratón) la carpeta con el tema del ratón que he almacenado
    previamente. Una vez hecho esto, actualizo con \texttt{update-alternatives}
    la referencia para que apunte al archivo \texttt{cursor.theme} del tema que
    he copiado.

    Una vez hecho eso, al reiniciar el sistema se actualizará el tema del ratón.

    \section{Conexión por ssh a mi cuenta de GitHub}
    Seguir las instrucciones de GitHub para conseguirlo.

    \section{Instalación de Syncthing}
    Seguir las indicaciones de la página web oficial de Syncthing.

    También tengo que copiar exactamente igual la carpeta que se especifique en
    el manual de Syncthing en la que guarda la configuración, que actualmente es
    \texttt{$\sim$/.local/state/syncthing/} y copiarla directamente antes de
    arrancar Syncthing (y después de instalarlo). La carpeta \texttt{index} que
    hay dentro de la carpeta con la configuración de Syncthing es mejor no
    copiarla porque es posible que el formato de almacenar esa información
    cambia si instalo Synchting de nuevo (aunque sea la misma versión porque
    mantenga actualizada la del ordenador), por tanto solamente copio el resto.

    \section{Instalación de Rclone}
    Instalar Rclone siguiendo las instrucciones de instalación de su página web
    y copiar el archivo de configuración guardado donde Rclone tenga su archivo
    de configuración.

    Para los servicios de copia de seguridad automática de las carpetas
    Documentos y Escritorio, tengo que copiar los archivos \texttt{.service} en
    la carpeta correspondiente a los servicios del usuario (con
    \texttt{man 5 systemd.unit} actualmente aparecen en las rutas en las que
    busca).

    \section{Instalación de Qterminal}
    No hace falta más configuración que la de establecer la fuente a la que
    tenga instalada para los iconos de \texttt{lf}, que reinstalaré más adelante.

    \section{Instalación de \LaTeX{}}
    Simplemente sigo las instrucciones de la página web para instalar la
    distribución de \LaTeX{} que prefiera.

    \subsection{Migración de mis paquetes}
    Tengo que crear una carpeta \texttt{texmf}, de acuerdo con las directrices
    actuales para añadir paquetes propios a LaTeX con los paquetes que tenga en
    mi anterior ordenador.

    \section{Instalación de Vimtex completa}
    Seguir las instrucciones del archivo correspondiente. Nótese que al copiar
    la carpeta \texttt{.vim} y el \texttt{.vimrc} hay que indicar en la
    ruta de configuración de vim el modo deseado del \texttt{.vimrc} del
    usuario, que será el de la aplicación del \texttt{.vimrc} por defecto
    y el del usuario, pues el \texttt{.vimrc} que guardo está reducido a
    solamente las configuraciones que guardo.

    Actualmente el archivo en el que se define ese comportamiento está en
    \texttt{/etc/vim}.

    \section{Instalación de KeePassXC}
    Viene ya instalado en la carpeta Escritorio, pero tengo que crear un enlace
    simbólico en \texttt{bin} para poderlo iniciar desde el \texttt{dpanel}.
    También puedo actualizarlo simplemente descargando el último
    \texttt{AppImage} de su página oficial. Y modificando el enlace simbólico
    de KeePassXC para que apunte a la nueva versión.

    \section{Instalación del tema de las aplicaciones GTK}
    El tema que tengo instalado está descargado de
    \href{https://github.com/vinceliuice/Graphite-gtk-theme}
    {https://github.com/vinceliuice/Graphite-gtk-theme}. En la configuración de
    apariencia de LXQt hay que habilitar los temas para GTK, pues si no no
    funcionará correctamente en las aplicaciones que lo soporten.
    
    \section{Instalación del tema de GRUB}
    La página para instalar el tema para Thunar también contendrá información
    sobre la instalación del mismo tema en GRUB.

    \section{Instalación del correo de Tuta}
    Simplemente instalo su aplicación e inicio sesión. Al igual que con
    KeePassXC (consultar su manual), hay que hacer un apaño para poder
    usarlo con \texttt{dmenu}.

    \section{Instalación del BlueTooth}
    Instalar \texttt{blueman}.

    \section{Instalación de Tor}
    Seguir las instrucciones de instalación para instalarlo en \texttt{opt}.

    \section{Instalación de lf}
    Seguir las instrucciones de su manual.

    \section{Instalación de DWM}
    La aplicación para realizar las capturas de pantalla será
    \texttt{screengrab} (o la aplicación que venga por defecto para realizar
    capturas de pantalla con el entorno de escritorio). En caso de cambiar
    cualquier otra aplicación también tendré que cambiarla debidamente en el
    archivo de configuración de DWM.

    En el caso de los archivos de configuración es mejor hacerlo como en el caso
    de SearxNG, comparar el archivo de configuración actual con la última
    versión que haya, y modificar el de la última versión en consecuencia. Hay
    que tener cuidado con modificar también la tecla \texttt{MODKEY}.

    Antes de iniciar dwm hay que comprobar que el guion que genera la barra
    funcione correctamente, pues está mal hecho y puede fallar dependiendo del
    sistema.

    Ahora basta decirle a LXQt el tema que debe poner (para el gestor de
    archivos y el resto de aplicaciones del sistema que siga usando) en DWM.
    Para ello, debo consultar en la wiki de LXQt dónde se encuentran los
    archivos de configuración. Los locales actualmente se encuentran en
    \texttt{$\sim$/.config/lxqt}, y los globales en \texttt{/etc/xdg/lxqt}.
    Tengo que copiar los locales (una vez configurado los ajustes que se deseen
    dentro de LXQt usando su aplicación gráfica para establecerlos. Véase
    \href{https://wiki.archlinux.org/title/Qt#Configuration_of_Qt_5/6_applications_under_environments_other_than_KDE_Plasma}
    {https://wiki.archlinux.org/title/Qt\#Configuration\_of\_Qt\_5/6\_applications\_under\_\-environments\_\-other\_than\_KDE\_Plasma}.

    En el archivador que viene por defecto hay que tener cuidado, pues las
    opciones de comprimir y descomprimir no funcionan, pues por defecto en la
    configuración del gestor de archivos aparece que la herramienta para los
    archivos comprimidos es \texttt{xarchiver}, cuando debería ser
    \texttt{lxqt-archiver}.

    Usando actualmente \texttt{lxappearance} se puede configurar cómo se ven
    en general las aplicaciones. Para arreglar que el ratón no se vea sin el
    tema instalado en los bordes de DWM hay que seguir las instrucciones del
    enlace
    \href{https://unix.stackexchange.com/questions/541766/change-icon-and-cursor-theme-in-dwm}
    {https://unix.stackexchange.com/questions/541766/change-icon-and-cursor-theme-in-dwm}.

    Ahora consulta el manual de la configuración de gestor de ventanas para más
    información.

    Es posible que para que se mantengan los ajustes de apariencia de las
    aplicaciones sea necesario tener instalado el paquete
    \texttt{lxde-settings-daemon}.

    \section{Utilización de NextDNS}
    Iniciar sesión en Internet en la cuenta de NextDNS para el ordenador y configurarla
    tanto en el navegador como en el sistema siguiendo las instrucciones que indican.

    \section{Terminación}
    Desinstalar las herramientas como \texttt{wget} y \texttt{curl} que se hayan
    necesitado instalar para instalar algunas de las aplicaciones. Y comprobar
    que Rclone esté funcionando correctamente.

    \section{Solución de errores}

    \subsection{Error en WiFi}
    El error en el WiFi se soluciona usando el comando \texttt{nmcli}. Actualmente
    lo que hay que hacer es:

    \begin{minted}[breaklines, breakanywhere, breaksymbolleft="", breaksymbolright="", breakanywheresymbolpre="", breakanywheresymbolpost=""]{bash}
nmcli connection show # Me muestra las redes que hay guardadas, copio el UUID
# de la red a la que me quiero conectar.
nmcli connection edit UUID
set connection.autoconnect=yes
set connection.autoconnect-retries=0    # Intenta siempre reconectarse.
    \end{minted}
\end{document}
